\documentclass[11pt]{letter}
\usepackage[T1]{fontenc}
\usepackage{lmodern}
\usepackage[margin=2.5cm]{geometry}
\usepackage{xurl}
\usepackage[hidelinks]{hyperref}
\signature{Minas Aslanyan\\Institute for Informatics and Automation Problems\\National Academy of Sciences of the Republic of Armenia, Yerevan\\\href{mailto:minas.aslanyan@iiap.sci.am}{minas.aslanyan@iiap.sci.am}}
\address{}
\date{\today}

\begin{document}
\begin{letter}{The Editors-in-Chief\\Pattern Recognition}
\opening{Dear Editors,}

We submit the manuscript ``Training parts versus matching parts in pose-based
gait recognition'' for
consideration as a regular research article in \emph{Pattern Recognition}.

Part-aware gait recognizers usually add two things at once: identity losses on
body regions during training and fusion of regional similarities at
retrieval. Gains are then credited to one of them without testing the other.
Our study separates the two. On one fixed pose-graph encoder, seven training
objectives, including global and part supervised contrastive losses and
cross-entropy controls, were trained with five seeds each on CASIA-B
development identities. No auxiliary objective differed reliably from a
metric-learning baseline, and an apparent penalty of part supervision in an
earlier three-seed experiment did not replicate. In contrast, fusing normalized
raw part pools from frozen checkpoints improved corrupted-probe accuracy under
every objective. On Gait3D, with the recipe frozen on development identities
and the test identities evaluated once, the same matcher raised Rank-1 from
21.88\% to 23.86\% and reached parity with GPGait, which we reproduced with its
official configuration on identical data.

We then applied the matcher to GPGait itself, and the result reversed: GPGait's
part-specific branches, trained with their own losses, add 9.1 points on
CASIA-B and 5.2 on Gait3D, while raw part matching lowers its accuracy. The
paper therefore concludes that where a part-aware gain comes from depends on
the encoder, and that global, raw regional and trained regional matching should
be compared on the same frozen checkpoints. We report the limitations openly,
including a large accuracy gap to GPGait on CASIA-B, earlier project exposure to
the CASIA-B test identities, an available-pose subset of Gait3D and the absence
of labelled real occlusion.

We believe the paper suits \emph{Pattern Recognition} because it addresses a
methodological question about how part-aware recognition systems are evaluated,
and supports its answer with controlled experiments, seed-by-identity bootstrap
intervals, reproduced baselines and a transfer test on an established
encoder. Every table and
figure can be regenerated from the saved per-probe outcomes with the released
scripts.

This manuscript is original, has not been published, and is not under
consideration for publication elsewhere. Both authors approved the submission
and declare no competing interests. The study received no specific funding.
The use of an AI assistant during manuscript preparation is declared in the
manuscript.

\closing{Sincerely,}
\end{letter}
\end{document}
